\documentclass[12pt]{article}

\usepackage{amsmath}
\DeclareMathOperator{\tr}{tr} % or \tr / \trace

\usepackage{bbm}
\usepackage{fancyhdr}
\usepackage{titling}
\usepackage{enumitem}
\usepackage{amsmath}
\usepackage{dsfont}
\usepackage{amsfonts}
\usepackage{amssymb}
\usepackage{amsthm}
\usepackage{geometry}
\usepackage{tikz}
\usepackage{algorithm}
\usepackage{algpseudocode}
\geometry{verbose,tmargin=1in,bmargin=1in,lmargin=1in,rmargin=1in}
\renewcommand{\vec}[1]{\mathbf{#1}}
\newcommand{\dec}{\text{dec}}
\newcommand{\divi}{\mathbin{\textnormal{div}}}
\title{Homework 1}
\author{Satvik Prasad}
\date{}
\newtheorem{lemma}{Lemma}[section]
\fancypagestyle{plain}{%
  \fancyhf{}
  \fancyhead[L]{\theauthor}
  \fancyhead[R]{\thedate}
}
\pagestyle{plain}
\makeatletter
\def\@maketitle{%
  \newpage
  \null
  \vskip 1em%
  \begin{center}%
    \let \footnote \thanks
    {\LARGE \@title \par}%
    \vskip 1em%
  \end{center}%
  \par
\vskip 1em}
\makeatother
\usepackage{cmbright}

\newcommand{\mat}[1]{\mathbf{#1}}

% ---------- Probability theory ----------
% Core operators (optional argument-free; use as \E{X}, \Var{X}, \Cov{X}{Y})
\DeclareMathOperator{\EOp}{\mathbb{E}}
\DeclareMathOperator{\VarOp}{Var}
\DeclareMathOperator{\CovOp}{Cov}
\DeclareMathOperator{\CorrOp}{Corr}

\renewcommand{\P}[1]{\mathbb{P}\left(#1\right)}        % \P{A}
\newcommand{\E}[1]{\EOp\left[#1\right]}                % \E{X}
\newcommand{\Var}[1]{\VarOp\left(#1\right)}            % \Var{X}
\newcommand{\Cov}[2]{\CovOp\left(#1,\,#2\right)}       % \Cov{X}{Y}
\newcommand{\Corr}[2]{\CorrOp\left(#1,\,#2\right)}     % \Corr{X}{Y}
\newcommand{\SD}[1]{\operatorname{SD}\left(#1\right)}  % \SD{X}

% Subscripted versions, e.g. \Esub{\theta}{X}, \Psub{\theta}{A}
\newcommand{\Esub}[2]{\EOp_{#1}\left[#2\right]}
\newcommand{\Psub}[2]{\mathbb{P}_{#1}\left(#2\right)}

% Conditioning, indicators, independence
\newcommand{\given}{\,\middle|\,}                      % \E{X \given Y}
\newcommand{\ind}[1]{\mathbbm{1}\left\{#1\right\}}     % \ind{X > 0}
\newcommand{\indep}{\perp\!\!\!\perp}                  % X \indep Y
\newcommand{\iid}{\overset{\text{iid}}{\sim}}          % X_i \iid \Normal(0,1)

% Convergence
\newcommand{\convp}{\xrightarrow{\;p\;}}               % in probability
\newcommand{\convd}{\xrightarrow{\;d\;}}               % in distribution
\newcommand{\convas}{\xrightarrow{\text{a.s.}}}        % almost surely
\newcommand{\convLp}[1]{\xrightarrow{L^{#1}}}          % in L^p

% Probability spaces
\newcommand{\F}{\mathcal{F}}                           % sigma-algebra
\newcommand{\B}{\mathcal{B}}                           % Borel sigma-algebra
\newcommand{\pspace}{(\Omega, \F, \mathbb{P})}

% Common distributions
\newcommand{\Normal}{\mathcal{N}}
\newcommand{\Bern}{\operatorname{Bernoulli}}
\newcommand{\Bin}{\operatorname{Binomial}}
\newcommand{\Pois}{\operatorname{Poisson}}
\newcommand{\Unif}{\operatorname{Uniform}}
\newcommand{\Exp}{\operatorname{Exp}}
\newcommand{\Geom}{\operatorname{Geometric}}
\newcommand{\Gam}{\operatorname{Gamma}}
\newcommand{\Beta}{\operatorname{Beta}}

% Number sets
\newcommand{\R}{\mathbb{R}}
\newcommand{\N}{\mathbb{N}}
\newcommand{\Z}{\mathbb{Z}}


\begin{document}
\subsection{Green Circle}
\begin{enumerate}
    \item Assume that $A^3 = 0$ for some real square matrix. Then, $I - A^3 = I.$ Equivalently, factorizing yields
    $$
    (I - A)(A^2 + A + I) = I,
    $$
    meaning $A^2 + A + I$ is a right inverse of $I - A.$

    Therefore, since $A$ is square we know $A^2 + A + I$ must also be a left inverse of $I - A,$ meaning $I - A$ is invertible.
    
    \item[4] Assume we have $A, B \in \mathbb R^{n \times n}$ satisfying $$
    \tr{A A^T + B B^T} = \tr{AB + A^T B^T}.
    $$

    Then, we know \begin{align*}
        \tr({A A^T + B B^T - AB - A^T B^T})  = 0.
    \end{align*}

    Equivalently, we get 
    $$
    \tr({A(A^T - B) + (B - A^T)B^T}) = 0.
    $$
    Since $\tr AB = \tr BA,$ we can commute inside the trace operator, meaning we have
    $$
    \tr((A - B^T)(A^T - B)) = 0.
    $$
    Let $C = A- B^T.$ We want to prove that $C = 0,$ given $\tr C C^T = 0.$

    Letting $C = \begin{pmatrix}
        v_1 & v_2 & \ldots & v_n
    \end{pmatrix}$
    for $v_i \in \mathbb R^{n},$ we know the diagonal entries of $C C^T$ are $||v_i||^2.$ Hence, $$\tr C C^T = \sum_{i = 0}^{n} ||v_i||^2 = 0.$$
    Since each term is $\geq 0,$ this forces $||v_i|| = 0,$ which requires $v_i = 0$ for all $1 \leq i \leq n,$ meaning $C = 0.$ Thus, we have $A - B^T = 0,$ meaning $A = B^T.$

    \item[5] Since very single column of $A$ has column-sum $m,$ we can know the following: $$
    A_{nj} = m - \sum_{i = 1}^{n -1 } A_{ij}
    $$ 
    for all $1 \leq j \leq n.$ Consider the map $T : \mathbb R^{n \times n} \to \mathbb R^{n \times n}$  given by 
    $$
    T = \begin{pmatrix}
        1 & 0 & \ldots & 0 & 0 \\  
        0 & 1 & \ddots & 0 & 0 \\
        0 & 0 & \ldots & 1 & 0 \\
        1 & 1 & \ldots & 1 & 1
    \end{pmatrix}
    $$
    since $\det T = 1$ (lower triangular), we know $$
    \det A = \det TA = \det A^\star,
    $$
    where $A^\star_{i \cdot} = A_{i \cdot}$ for all $1 \leq i < n,$ and $A^\star_{nj} = m$ for all $1 \leq j \leq n.$ Then, we have
    $$
    \det A^\star = \sum_{\sigma \in S_n} \text{sign}(\sigma) \prod_{i = 1}^{n} A^\star_{i \sigma(i)} = m \sum_{\sigma \in S_n} \text{sign}(\sigma) \prod_{i = 1 }^{n - 1} A_{i \sigma(i)}.
    $$
    Since $A$ has integer entries, we can conclude that $$
    \det A =\det A^\star \equiv 0 \pmod {m},
    $$
    meaning the determinant of $A$ is divisible by $m.$
\end{enumerate}

\subsection{Blue Square}
\begin{enumerate}
    \item[6] Equivalently, we have $I + A^2 = 0,$ where $I$ is the $n \times n$ identity matrix. Let us analyze the eigenvalues of $A.$ For some arbitrary vector $v \neq 0,$ $$
    Av =  \lambda v \Rightarrow A^2v = \lambda ^2 Iv =  -Iv,
    $$
    meaning $(\lambda^2 + 1)I v = 0,$ or equivalently, $\lambda = \pm i.$ Therefore, $A$ has two eigenvalues, $i$ and $-i.$ 

    Hence, the characteristic polynomial $p \in \mathbb R[\lambda]$ produced by $\det (A - \lambda I) = 0$ only has roots $i$ and $-i,$ and satisfies $$p(0) = 0.$$ By Vieta's, $\alpha i - \beta i = (\alpha - \beta)i$  must be in $\mathbb R,$ meaning $\alpha = \beta.$ Hence, we can write the characteristic polynomial in the form
    $$
    p(\lambda) = c(\lambda - i)^{\alpha}(\lambda + i)^{\alpha}= c(\lambda^2 + 1)^{\alpha},
    $$
    and thus $p(0) = c.$ Since the characteristic polynomial is always monic, we know $c = 1,$ and thus $\det A = p(0) = 1.$

    \item[8] We have square matrices $A, B.$ We proceed by showing that 
    $$
    \dim \ker (I - AB) > 0 \Leftrightarrow \dim \ker (I - BA) > 0.
    $$

    If $\dim \ker (I - AB) > 0,$ there is some non-zero $v$ such that
    $$
    v - ABv = 0 \Rightarrow ABv = v.
    $$

    Multiplying both sides by $B$, we get $BABv = Bv,$ meaning $BAu = u,$ where $u = Bv.$ Therefore, 
    $$
    (I - BA) u = 0.
    $$
    If $v \in \ker B,$ then $v - ABv = v = 0,$ implying $v = 0.$ Therefore, $\ker B = \{0\}.$ Since $v \neq 0,$ $u = Bv \neq 0,$ meaning $\dim \ker (I - BA) > 0.$

    Conversely, if $\dim \ker (I - BA) > 0,$ then there is some $v \neq 0$ such that
    $$
    v - BAv = 0 \Rightarrow BAv = v.
    $$
    Then, multiplying both sides by $A,$ we get $$
    AB(Av) =Av \Rightarrow AB u = u,
    $$
    where $u = Bv.$ By a similar argument, $\ker A = \{0\},$ meaning $u \neq 0,$ and hence $\dim \ker (I - AB) > 0.$

    By contraposition,
    $$
    \dim \ker (I - AB) = 0 \Leftrightarrow \dim \ker (I - BA) = 0,
    $$
    meaning $I - AB$ is invertible iff $I - BA$ is invertible.

    \item[9] We have real numbers $a, b$ such that $D_n(a, b)$ be the matrix with diagonal entries $a$ and all other entries $b.$ When $b = 0,$ we know $D_{n}(a, b) = a^{n}.$

    Define the matrix $T \in \mathbb R^{n \times n}$ such that
    $$
    T = \begin{pmatrix}
        \frac{a + (n-2)b} {b} & -1 & \ldots & -1  & -1 \\
        0 & 1 &  \ldots & 0 & 0 \\
        \vdots & \vdots & \ddots &\vdots  & \vdots \\ 
        0 & 0 & \ldots & 1 & 1 \\
        0 & 0 & \ldots & 0 & 1 \\
    \end{pmatrix},
    $$
    where $\det T = (a + (n-2)b) / b$ since $T$ is upper-triangular. Then $T D_n(a, b)$ satisfies
    \begin{align*}
    TD_{n}(a, b)_{1 \cdot }&= \dfrac{a + (n-2)b}{b} \begin{pmatrix}
        a \\
        b \\ 
        \vdots \\
        b 
    \end{pmatrix}  - \begin{pmatrix}
        b \\
        a  \\ 
        \vdots \\
        b
    \end{pmatrix} - \ldots - \begin{pmatrix}
        b \\
        b \\
        \vdots \\
        a
    \end{pmatrix} \\
    &= \left[\dfrac{a^2}{b} + (n-2)a - (n-1)b\right]\begin{pmatrix}
        1 \\
        0 \\
        \vdots \\
        0
    \end{pmatrix}
    \end{align*}

    Therefore, applying cofactor expansion yields
    $$
    \det T D_{n}(a, b) = \left[\dfrac{a^2}{b} + (n-2)a - (n-1)b\right]  \det  D_{n-1}(a, b),
    $$
    where $D_1(a, b) = a.$ By the properties of the determinant, we get
    $$
    \det D_{n}(a, b) = \dfrac{1}{a + (n-2)b} \left[a^2 + (n-2)ab - (n-1)b^2\right] \det D_{n-1}(a, b).
    $$

    Now, let us prove by induction that $$
    \det D_n(a, b) = (a-b)^{n-1}(a + (n-1)b) := K(n).
    $$

    For $n = 1,$ $D_n(a, b) = a,$ and $K(1) = a$ as required. Assume for $n = k \geq 1,$ we know $$
        \det D_k(a, b) = K(k) = (a-b)^{k-1}(a + (k-1)b).
    $$
    Then, $D_{k + 1}(a, b)$ satisfies \begin{align*}
            \det D_{k+1}(a, b) &= \left[a - \dfrac{kb^2}{a + (k-1)b} \right] (a-b)^{k-1} (a + (k-1)b) \\
            &=\left[a(a + (k -1)b) - kb^2\right] (a-b)^{k-1} \\
            &= (a^2 - kb^2 + (k-1)ab) (a-b)^{k-1} \\
            &= (a^2 - ab + kab - kb^2) (a-b)^{k-1} \\
            &= (a(a- b) + kb(a - b)) (a-b)^{k-1} \\
            &= (a-b)^k (a + (k + 1 -1)b) = K(k + 1)
    \end{align*}
    as required.

    Therefore, by induction we know
    $$
    \det D_n(a, b) = (a-b)^{n-1} (a + (n - 1)b)
    $$
    for all integers $n \geq 1$ as required.
\end{enumerate}

\subsection{Black Diamond}
\begin{enumerate}
    \item[12] Define $\sin A$ such that
    $$
    \sin A = \sum_{k = 0}^{\infty} \dfrac{(-1)^{k}}{(2k+1)!} A^{2k + 1}.
    $$
    We want to prove that there is no such matrix for which $$
    \sin A = \begin{pmatrix}
    1 & 2027 \\
    0 & 1
    \end{pmatrix}.
    $$

    For the sake of contradiction, assume the contrary, i.e., there is some $A$ such that $\sin A = \square$.

    If $A$ has an eigenvalue $\lambda,$ $Av = \lambda v$ for some non-zero $v.$ Then,
    $$
    (\sin A ) v = \sin(\lambda)v,
    $$
    meaning $\sin \lambda $ is an eigenvalue of $\sin A.$
    
    Assume $\sin A$ has some eigenvalue $\mu \in \mathbb R.$ Then, for some $v \neq 0,$ $$
    \sum_{k = 0}^{\infty} \dfrac{(-1)^k}{(2k+1)!} A^{2k+1} v = \mu.
    $$

    Then, we have $$
    \dim \ker \left(\sin A- \mu I\right) > 0.
    $$

    For some $v \in \ker (\sin A - I),$ we know $$(\sin A)  \cdot A v = A \cdot  (\sin A) v = A (0) = 0,$$
    meaning $Av  \in \ker(\sin A - \mu I),$ and hence $\ker (\sin A - \mu I)$ is invariant under $A.$ Then, consider some $u \in \ker (\sin A - \mu I)$ that is non-zero. This means
    $$
        (\sin A ) \cdot u = \mu  u.
    $$

    Since $\dim \ker (\sin A - \mu I) >0,$ this subspace is either $1$-dimensional, or $2$-dimensional.

    \begin{enumerate}
        \item If $1$-dimensional, then $Au = \lambda u$ for some scalar $\lambda.$ Substituting, this yields $\sin \lambda \cdot u = \mu u,$ meaning $\mu = \sin \lambda.$ Hence, there is some eigenvalue $\lambda$ of $A$ such that $\sin \lambda = \mu.$

        \item Otherwise, $\text{range}(\sin A - \mu I) = \{0\}$ by rank-nullity, meaning $\sin A = \mu I.$ However, our specific $\sin A$ is not a scalar multiple of the identity, meaning this is false.
    \end{enumerate}

    Therefore, any eigenvalue $\lambda$ of $A$ must satisfy $\sin \lambda = 1.$

    Now, consider
    $$
        \cos A : = \sum_{k = 0}^{\infty} \dfrac{(-1)^k}{(2k)!} A^{2k}.
    $$
    If $A$ has an eigenvalue $\lambda,$ then $\cos A$ has eigenvalues $\cos \lambda$ by a similar argument to the above. Hence, $0$ is an eigenvalue of $\cos A.$

    If $\mu$ is an eigenvalue of $\cos A,$ then there must be some eigenvalue $\lambda$ of $A$ such that $\cos \lambda  = \mu$ by an equivalent approach over $\cos \cdot.$ Since all eigenvalues of $A$ satisfy $\sin \lambda = 1,$ all eigenvalues of $A$ must satisfy $\cos \lambda = 0,$ meaning $0$ is the only eigenvalue of $\cos A.$

    Hence, the characteristic polynomial of $\cos A$ is $p(\mu) = \mu^2,$ meaning by Cayley-Hamilton, $(\cos A)^2 = 0.$

    Looking at $(\sin A^2 + \cos A)^2,$ we have
    \begin{align*}
    (\sin A)^2 + (\cos A)^2 &= \sum_{k = 0}^{\infty} \sum_{j = 0}^{\infty} \left(\dfrac{(-1)^{k +j}}{(2k)! (2j)!} A^{2(k  + j)} + \dfrac{(-1)^{k + j}}{(2k+1)! (2j + 1)!} A^{2(k + j + 1)} \right) \\
    &= \sum_{i = 0}^{\infty} \sum_{j = 0}^{i} \dfrac{(-1)^i}{(2(i - j))!(2j)!} A^{2i} + \sum_{i = 1}^{\infty} \sum_{j = 0}^{i - 1} \dfrac{(-1)^{i - 1}}{((2(i - 1 - j) + 1)!)(2j+1)!} A^{2i} \\
    &= I + \sum_{i = 1}^{\infty} \left(\dfrac{(-1)^i}{(2i)!} + \sum_{j = 0}^{i - 1} \left\{\dfrac{(-1)^i}{(2(i-j))! (2j)!} + \dfrac{(-1)^{i - 1}}{(2(i - 1 - j) + 1)!(2j+1)!}\right\} \right)A^{2i} \\
    &= I + \sum_{i = 1}^{\infty} \left(\dfrac{(-1)^i}{(2i)!} + \sum_{j = 0}^{i - 1} \left\{\dfrac{(-1)^i(2j+1) + (-1)^{i-1}(2(i -j))}{(2(i -j))!(2j + 1)!}\right\} \right)A^{2i} \\    
    &= I + \sum_{i = 1}^{\infty} \left(\dfrac{(-1)^i}{(2i)!} + \sum_{j = 0}^{i - 1} \left\{\dfrac{(-1)^{i - 1}[2i - 4j - 1]}{(2(i -j))!(2j + 1)!}\right\} \right)A^{2i} \\
    &= I + \sum_{i = 1}^{\infty} \left(\dfrac{(-1)^i}{(2i)!} + \sum_{j = 0}^{i - 1} (-1)^{i - 1} \dfrac{2i - 4j - 1}{(2i+1)!}{2i + 1 \choose 2j + 1} \right)A^{2i}.
    \end{align*}

    Since the binomial coefficient is symmetric, \begin{align*}
    &\sin(A)^2 + \cos(A)^2 \\&= I + \sum_{i = 1}^{\infty} \left(\dfrac{(-1)^i}{(2i)!} + \dfrac{1}{(2i + 1)!}\sum_{j = 0}^{i - 1} (-1)^{i - 1} (2i - 2j - (2j + 1)){2i + 1 \choose 2j + 1} \right)A^{2i}.\\
    &= I + \sum_{i = 1}^{\infty} \left(\dfrac{(-1)^i}{(2i)!} + \dfrac{(-1)^{i - 1}}{(2i + 1)!}\left\{\sum_{j = 0}^{i - 1}(2i - 2j){2i + 1 \choose 2j + 1}  - \sum_{j = 0}^{i - 1} (2i - 2j - 1){2i + 1 \choose 2j + 2}\right\}\right)A^{2i}.
    \end{align*}
    Since $(n-r){n \choose r} = n {n -1 \choose r},$ we get \begin{align*}
    &\sin (A)^2 + \cos (A)^2 \\
    &= I + \sum_{i = 1}^{\infty} \left(\dfrac{(-1)^i}{(2i)!} + \dfrac{(-1)^{i - 1}}{(2i + 1)!}\left\{\sum_{j = 0}^{i - 1}(2i + 1){2i\choose 2j + 1}  - \sum_{j = 0}^{i - 1} (2i + 1){2i \choose 2j + 2}\right\}\right)A^{2i} \\
    &= I + \sum_{i=1}^{\infty}
\left(
  \dfrac{(-1)^i}{(2i)!}
  + \dfrac{(-1)^{i-1}}{(2i)!}
    \left[
      \sum_{j=0}^{i-1} \binom{2i}{2j+1}
      - \sum_{j=0}^{i-1} \binom{2i}{2j+2}
    \right]
\right) A^{2i} \\
&= I + \sum_{i=1}^{\infty}
\left(
  \dfrac{(-1)^i}{(2i)!}
  + \dfrac{(-1)^{i}}{(2i)!}
    \left[
      \sum_{k = 1}^{2i} {2i \choose k} (-1)^{k}
    \right]
\right) A^{2i}  \\
&= I + \sum_{i=1}^{\infty}
\left(
  \dfrac{(-1)^i}{(2i)!}
  + \dfrac{(-1)^{i}}{(2i)!}
    \left[
      (1 - 1)^{2i} - 1
    \right]
\right) A^{2i}  \\
&= I,
    \end{align*}
    as required.

    Therefore, $\sin(A)^2 = I - \cos(A)^2 = I,$ which is not the case, yielding a contradiction. \qed
\end{enumerate}

\subsection{Reflection}
Looked at eigenvalues for 12. The given matrix has only the eigenvalue \(1\), so every eigenvalue \(\lambda\) of \(A\) satisfies \(\sin\lambda=1\). This means both eigenvalues of \(\cos A\) are zero. The \(2\times2\) size mattered here because the characteristic polynomial of \(\cos A\) is \(t^2\), so Cayley–Hamilton gives \((\cos A)^2=0\).

The identity \((\sin A)^2+(\cos A)^2=I\) seemed intuitive and proveable, so brute-forced proved it. Squaring the given matrix produces an upper-right entry of \(4054\), giving the contradiction.
\end{document}